\chapter{METHODOLOGY}

This chapter details the methodology, mathematical foundations, algorithmic pipelines, and architectural design of the proposed mobile cardiovascular monitoring and digital-twin simulation platform. The system acquires non-invasive sternal seismocardiography (SCG) directly via consumer smartphone MEMS accelerometers, followed by automated digital signal conditioning, statistical pulse rate variability (PRV) extraction, supervised machine learning classification, and in-silico nonlinear dynamical rhythm simulation. To ensure reproducibility and transparency, this chapter is organized into five core sections: System Design, Data, Methods, Hyperparameter Configurations, and Performance Evaluation.

\section{System Design}

\subsection{Architectural Overview}
The platform is designed around a decoupled, four-tier pipeline. Each tier operates with standardized interface contracts so that signal degradation or motion artifacts in the sternal SCG acquisition stream are deterministically intercepted by physiological energy and SNR quality gates, triggering immediate real-time stillness and placement guidance before corrupted data can propagate downstream. Figure~\ref{fig:system_architecture} illustrates the high-level system architecture across its four functional tiers:

\begin{figure}[!htbp]
	\centering
	\fbox{\parbox{0.92\textwidth}{
			\centering
			\vspace{4pt}
			\textbf{TIER 1: Data Acquisition Layer}\par\vspace{3pt}
			{[Sternal MEMS Accelerometer ($\pm 2g$, 100--200 Hz)] $\longrightarrow$ [Real-Time Stillness Guidance]}\par\vspace{5pt}
			$\Downarrow$\par\vspace{5pt}
			\textbf{TIER 2: Signal Conditioning \& Pre-processing Layer}\par\vspace{3pt}
			{[Zero-Phase Butterworth Filter (0.5--40 Hz)] $\longrightarrow$ [Energy Gate: $0.005 \le \sigma_{raw} \le 0.5$ m/s$^2$] $\longrightarrow$ [SNR Check $\ge 12$ dB]}\par\vspace{5pt}
			$\Downarrow$\par\vspace{5pt}
			\textbf{TIER 3: Feature Extraction Layer}\par\vspace{3pt}
			{[Adaptive Sliding Threshold $\theta[n]$] $\longrightarrow$ [AO Peak Detection] $\longrightarrow$ [Tachogram Interpolation (4 Hz)] $\longrightarrow$ [PRV Metrics]}\par\vspace{5pt}
			$\Downarrow$\par\vspace{5pt}
			\textbf{TIER 4: Analytics \& In-Silico Simulation Layer}\par\vspace{3pt}
			{[ML Classifiers: Normal vs. Heart Problem / Arrhythmia \& Stress] $\longleftrightarrow$ [Coupled Van der Pol Dynamical ODE Simulation]}\par\vspace{4pt}
	}}
	\caption{Four-Tier System Architectural Block Diagram}
	\label{fig:system_architecture}
\end{figure}

The four functional layers are structured as follows:
\begin{enumerate}
	\item \textbf{Tier 1 --- Data Acquisition Layer:} Captures raw physiological cardiac mechanical signals directly on consumer smartphones. Sternal SCG micro-vibrations are captured along the dorso-ventral ($Z$-axis) by the phone's internal 3-axis MEMS accelerometer ($\pm 2g$, $100\text{--}200\text{ Hz}$) placed flat on the subject's sternum during resting supine postures.
	\item \textbf{Tier 2 --- Signal Conditioning Layer:} Attenuates respiratory drift, high-frequency muscular tremor, and out-of-band noise using 4th-order zero-phase Butterworth bandpass filtering. Evaluates signal viability using an absolute physiological energy gate ($0.005\text{ m/s}^2 \le \sigma_{raw} \le 0.50\text{ m/s}^2$) and computes quantitative Signal-to-Noise Ratio (SNR).
	\item \textbf{Tier 3 --- Feature Extraction Layer:} Locates aortic opening (AO) systolic peaks using an adaptive moving threshold. Constructs peak-to-peak interval (PPI) tachograms, enforces physiological validity bounds ($333\text{--}1500\text{ ms}$), and calculates standardized time-domain and frequency-domain PRV metrics.
	\item \textbf{Tier 4 --- Analytics and In-Silico Simulation Layer:} Maps the PRV feature vector into supervised machine learning classifiers to distinguish normal sinus rhythm from cardiac abnormalities / arrhythmias, while concurrently driving a coupled Van der Pol nonlinear cardiac oscillator model to evaluate dynamic rhythm stability under simulated stress challenges.
\end{enumerate}

\subsection{Operational Decision Flowchart}
The operational runtime execution follows an automated decision-loop protocol:
\begin{enumerate}
	\item \textbf{Measurement Initialization:} The user initiates recording, buffering a 30 to 60-second window of accelerometer samples.
	\item \textbf{Posture \& Placement Verification:} The system evaluates gravity baseline orientation. If the subject is stationary and supine ($|a_z| \approx 9.8\text{ m/s}^2$), sternal SCG processing proceeds. If posture reorientation or excessive motion is detected, the system prompts the user to maintain a stationary supine posture.
	\item \textbf{Energy \& SNR Quality Gate:} Post-filtering, the signal energy and SNR are computed. If the signal energy satisfies physiological bounds ($0.005\text{ m/s}^2 \le \sigma_{raw} \le 0.50\text{ m/s}^2$) and $\text{SNR} \ge 12\text{ dB}$, sternal systolic peak detection is executed. If corrupted, the recording is rejected with immediate placement guidance.
	\item \textbf{Peak Extraction \& PRV Calculation:} Systolic peaks are identified, PPI outliers outside $333\text{--}1500\text{ ms}$ are rejected, and PRV features are computed.
	\item \textbf{Classification \& Simulation Output:} The extracted feature vector is classified by the trained ML models (Normal vs. Heart Problem), and the digital-twin simulation renders dynamic phase-plane trajectories on the user dashboard.
\end{enumerate}

\subsection{Algorithmic Pseudocode}
Algorithm~\ref{alg:scg_dsp} outlines the complete mathematical sequence for sternal SCG signal conditioning, adaptive peak detection, and PRV feature extraction.

\begin{table}[!htbp]
	\centering
	\caption{Algorithm 1: Sternal SCG Conditioning and PRV Feature Extraction}
	\label{alg:scg_dsp}
	\begin{tabular}{|p{13.5cm}|}
		\hline
		\textbf{Input:} Raw dorso-ventral acceleration $a_z[n]$, sampling frequency $f_s \in [100, 200]\text{ Hz}$, window length $W=60\text{ s}$.\\
		\textbf{Output:} Heart rate ($\text{HR}$), $\text{SDNN}$, $\text{RMSSD}$, $\text{pNN50}$, $\text{LF/HF}$, Signal Quality Index ($\text{SQI}$).\\
		\hline
		\textbf{1. Physiological Energy Gating:} \\
		\quad Compute baseline noise standard deviation: $\sigma_{raw} = \text{std}(a_z)$.\\
		\quad \textbf{if} $\sigma_{raw} < 0.005\text{ m/s}^2$ \textbf{or} $\sigma_{raw} > 0.50\text{ m/s}^2$ \textbf{then}\\
		\quad \quad \textbf{return} status = \texttt{REJECTED\_UNVIABLE\_ENERGY}, prompt user to remain still and adjust sternal placement.\\
		\quad \textbf{end if}\\
		\textbf{2. Zero-Phase Bandpass Filtering:} \\
		\quad Design 4th-order Butterworth bandpass filter coefficients ($f_l = 0.5\text{ Hz}, f_h = 40.0\text{ Hz}$).\\
		\quad Apply forward-backward zero-phase filtering: $\tilde{a}_z = \text{FiltFilt}(a_z, b, a)$.\\
		\textbf{3. Signal Envelope Extraction:} \\
		\quad Compute moving root-mean-square (RMS) energy envelope over window $K = \lfloor 0.2 \cdot f_s \rfloor$ samples:\\
		\quad \quad $E[n] = \sqrt{\frac{1}{K}\sum_{m=-K/2}^{K/2} \tilde{a}_z[n+m]^2}$\\
		\textbf{4. Adaptive Systolic Peak Detection:} \\
		\quad Compute moving local mean $\mu_E[n]$ and moving standard deviation $\sigma_E[n]$ over $2\text{ s}$ window.\\
		\quad Define adaptive threshold: $\theta[n] = \mu_E[n] + \alpha \cdot \sigma_E[n]$ (with $\alpha = 0.6$).\\
		\quad Identify candidate indices $p_k$ where $E[p_k] > \theta[p_k]$ and $E[p_k] > E[p_k \pm 1]$.\\
		\quad Enforce physiological refractory period: $|p_k - p_{k-1}|/f_s \ge \delta_{ref} = 0.35\text{ s}$ ($\le 171\text{ bpm}$).\\
		\textbf{5. Tachogram Construction and Outlier Clamping:} \\
		\quad Compute beat intervals: $\text{PPI}_k = \frac{p_{k+1} - p_k}{f_s} \times 1000\text{ ms}$.\\
		\quad Reject ectopic outliers: retain only $\text{PPI}_k \in [333, 1500]\text{ ms}$.\\
		\textbf{6. Statistical Time-Domain PRV Extraction:} \\
		\quad $\text{HR} = 60000 / \text{mean}(\text{PPI})$.\\
		\quad $\text{SDNN} = \sqrt{\frac{1}{M-1}\sum_{k=1}^M (\text{PPI}_k - \overline{\text{PPI}})^2}$.\\
		\quad $\text{RMSSD} = \sqrt{\frac{1}{M-1}\sum_{k=1}^{M-1} (\text{PPI}_{k+1} - \text{PPI}_k)^2}$.\\
		\quad $\text{pNN50} = \frac{\text{Count}(|\text{PPI}_{k+1} - \text{PPI}_k| > 50\text{ ms})}{M-1} \times 100\%$.\\
		\textbf{7. Spectral Frequency-Domain PRV Extraction:} \\
		\quad Interpolate non-equidistant $\text{PPI}$ series to uniform $4.0\text{ Hz}$ grid using cubic splines.\\
		\quad Compute Welch Power Spectral Density: $P(f) = \text{Welch}(\text{PPI}_{4Hz})$.\\
		\quad Integrate frequency bands: $\text{LF} = \int_{0.04}^{0.15} P(f)df$, $\text{HF} = \int_{0.15}^{0.40} P(f)df$.\\
		\quad Compute sympathovagal balance: $\text{LF/HF} = \text{LF} / \text{HF}$.\\
		\textbf{8. Return:} $\{\text{HR}, \text{SDNN}, \text{RMSSD}, \text{pNN50}, \text{LF/HF}, \text{SQI}\}$.\\
		\hline
	\end{tabular}
\end{table}

\section{Data}

\subsection{Collection Setup and Hardware Interfacing}
Primary cardiac data is acquired using standard smartphone sensors without external peripheral electronics:
\begin{itemize}
	\item \textbf{Sternal Seismocardiography (SCG):} The smartphone is positioned flat on the lower third of the subject's sternum (xiphoid process region) while the subject rests in a supine posture \cite{inan2015ballistocardiography, tadi2016realtime}. The phone's internal tri-axial MEMS accelerometer is sampled at $100\text{ Hz}$ to $200\text{ Hz}$ through native mobile APIs (\texttt{CoreMotion} on iOS, \texttt{SensorManager} on Android). Sternal recoil vibrations along the dorso-ventral axis ($a_z$) directly reflect ventricular ejection forces and aortic valve opening (AO) dynamics.
\end{itemize}

\subsection{Acquisition Windows and Posture Gating}
Each recording session captures a contiguous $30\text{ to }60\text{-second}$ window. The $60\text{-second}$ duration represents the standard clinical baseline for short-term pulse rate variability, ensuring adequate frequency resolution ($\Delta f \approx 0.016\text{ Hz}$) to separate the LF band ($0.04\text{--}0.15\text{ Hz}$) from the HF band ($0.15\text{--}0.40\text{ Hz}$) without fatiguing the user. Posture gating continuously verifies that static gravity orientation remains within $\pm 10^\circ$ of the vertical $Z$-axis.

\subsection{Signal Viability and Real-Time Stillness Guidance Protocol}
When sternal SCG signal quality degrades---detected either by the absolute physiological energy gate ($\sigma_{raw} < 0.005\text{ m/s}^2$ indicating stationary surface placement, or $\sigma_{raw} > 0.50\text{ m/s}^2$ indicating active subject movement) or when $\text{SNR} < 12\text{ dB}$---the system immediately flags the SCG channel as unviable and halts feature extraction. Instead of generating spurious measurements from contaminated motion data, the application displays immediate on-screen guidance prompting the subject to remain still, maintain a relaxed supine posture, and check flat sternal placement before re-initiating data acquisition.

\subsection{Clinical Reference Datasets}
To benchmark the signal processing engine and train supervised machine learning models against established clinical ground truth, two publicly accessible PhysioNet and scientific repositories are utilized:
\begin{enumerate}
	\item \textbf{PhysioNet MIT-BIH Arrhythmia Database \cite{goldberger2000physiobank}:} Contains 48 half-hour two-channel ambulatory ECG recordings from 47 subjects, annotated beat-by-beat by expert cardiologists. This database provides the gold-standard ground truth for training and testing our binary cardiac classifier (Normal Sinus Rhythm vs. Heart Problem / Arrhythmia, encompassing premature ventricular contractions [PVCs], premature atrial contractions [PACs], and rate dysrhythmias).
	\item \textbf{WESAD (Wearable Stress and Affect Detection) Dataset \cite{schmidt2018wesad}:} Contains continuous physiological recordings (chest-worn RespiBAN and wrist-worn Empatica E4) from 15 subjects experiencing controlled emotional and autonomic stress challenges. This dataset provides annotated ground truth for evaluating our autonomic stress classification tier.
\end{enumerate}

\section{Methods}

\subsection{Digital Signal Conditioning}
Raw accelerometer signals exhibit DC gravity offsets ($~9.8\text{ m/s}^2$), low-frequency respiratory baseline wander ($0.1\text{--}0.5\text{ Hz}$), and high-frequency muscular noise ($>50\text{ Hz}$). To isolate cardiac micro-vibrations without introducing phase delay that would distort peak timings, a digital 4th-order zero-phase Butterworth bandpass filter \cite{butterworth1930theory} is applied using forward-backward recursive filtering (\texttt{scipy.signal.filtfilt}):
\begin{equation}
	|H(j\omega)|^2 = \frac{1}{1 + \left(\frac{\omega^2 - \omega_l \omega_h}{\omega(\omega_h - \omega_l)}\right)^{2N}}
\end{equation}
where $N=4$ is the filter order, $\omega_l = 2\pi(0.5)\text{ rad/s}$ is the lower cutoff frequency, and $\omega_h = 2\pi(40.0)\text{ rad/s}$ is the upper cutoff frequency.

\subsection{Energy Envelope Extraction}
Because filtered SCG waveforms exhibit complex bi-phasic oscillations, identifying true systolic aortic opening (AO) events directly on the raw oscillatory waveform can cause double-counting. We construct a unipolar energy envelope $E[n]$ using a moving root-mean-square (RMS) operator over window $K = \lfloor 0.2 \cdot f_s \rfloor$ samples:
\begin{equation}
	E[n] = \sqrt{\frac{1}{K}\sum_{m=-\lfloor K/2 \rfloor}^{\lfloor K/2 \rfloor} \tilde{a}_z[n+m]^2}
\end{equation}
This operation converts oscillatory wave-packets into smooth, distinguishable unimodal energy peaks coinciding with ventricular contraction.

\subsection{Adaptive Threshold Systolic Peak Detection}
To handle inter-subject amplitude variability and sensor coupling differences, peak detection employs an adaptive sliding threshold $\theta[n]$:
\begin{equation}
	\theta[n] = \mu_E[n] + \alpha \cdot \sigma_E[n]
\end{equation}
where $\mu_E[n]$ and $\sigma_E[n]$ are the running local mean and standard deviation of $E[n]$ over a sliding 2.0-second window, and $\alpha = 0.6$ is an empirically optimized scaling coefficient. A sample index $p_k$ is classified as a valid cardiac peak if:
\begin{equation}
	E[p_k] > \theta[p_k] \quad \text{and} \quad E[p_k] > E[p_k \pm 1] \quad \text{and} \quad \frac{p_k - p_{k-1}}{f_s} \ge \delta_{ref}
\end{equation}
where $\delta_{ref} = 0.35\text{ s}$ is the physiological refractory period, corresponding to an upper heart rate ceiling of $171\text{ bpm}$.

\subsection{Pulse Rate Variability (PRV) Feature Extraction}
From the extracted systolic peak timestamps, peak-to-peak intervals are computed: $\text{PPI}_k = (p_{k+1} - p_k)/f_s \times 1000\text{ ms}$. Intervals outside $333\text{--}1500\text{ ms}$ are clamped. The following standardized time-domain and frequency-domain metrics are extracted in adherence with Task Force guidelines \cite{taskforce1996heart, tarvainen2014kubios, schafer2013accurate}:
\begin{itemize}
	\item \textbf{Mean Heart Rate (Mean HR):} $\text{HR} = 60000 / \overline{\text{PPI}}\text{ (bpm)}$.
	\item \textbf{SDNN:} Standard deviation of all valid PPI intervals, reflecting overall autonomic variability.
	\item \textbf{RMSSD:} Root mean square of successive interval differences, reflecting vagal parasympathetic tone:
	\begin{equation}
		\text{RMSSD} = \sqrt{\frac{1}{M-1}\sum_{k=1}^{M-1} (\text{PPI}_{k+1} - \text{PPI}_k)^2}
	\end{equation}
	\item \textbf{pNN50:} Percentage of successive intervals differing by more than $50\text{ ms}$.
	\item \textbf{Spectral Powers (LF, HF, LF/HF):} The discrete non-equidistant PPI tachogram is resampled to a uniform $4.0\text{ Hz}$ time grid via cubic-spline interpolation. Welch's periodogram \cite{welch1967use} computes the power spectral density $P(f)$. Low-frequency power ($\text{LF}: 0.04\text{--}0.15\text{ Hz}$) and high-frequency power ($\text{HF}: 0.15\text{--}0.40\text{ Hz}$) are integrated. The ratio $\text{LF/HF}$ quantifies sympathovagal autonomic balance.
\end{itemize}

\subsection{Supervised Machine Learning Classification}
Extracted PRV feature vectors $\mathbf{x} = [\text{HR}, \text{SDNN}, \text{RMSSD}, \text{pNN50}, \text{LF}, \text{HF}, \text{LF/HF}]$ are evaluated across four supervised machine learning algorithms using \texttt{scikit-learn} \cite{pedregosa2011scikit}: Decision Tree (DT), Logistic Regression (LR), Random Forest (RF) \cite{breiman2001random}, and Extreme Gradient Boosting (XGBoost) \cite{chen2016xgboost}. Effective cardiac arrhythmia classification requires rigorous feature extraction and selection to distinguish normal physiological patterns from pathological beat formations \cite{pandian2026internet}. The models are trained to perform two classification tasks:
\begin{enumerate}
	\item \textbf{Primary Classification --- Cardiac Health Status (Binary \& Multi-Class):}
	Following clinical cardiac arrhythmia (CA) frameworks and AAMI standards \cite{pandian2026internet}, cardiac beats are categorized into five fundamental clinical beat types:
	\begin{itemize}
		\item \textit{Normal (N):} Normal sinus beats and bundle branch block complexes.
		\item \textit{Supraventricular ectopic (S):} Atrial premature beats (PACs) and supraventricular arrhythmias.
		\item \textit{Ventricular ectopic (V):} Premature ventricular contractions (PVCs) and ventricular escape beats.
		\item \textit{Fusion (F):} Fusion of ventricular and normal conduction impulses.
		\item \textit{Unknown / Unclassifiable (Q):} Paced beats and unclassified dysrhythmias.
	\end{itemize}
	In our mobile screening prototype, these are evaluated both as an overarching binary screening model (Class 0: Normal [N] vs. Class 1: Heart Problem / Arrhythmia [S, V, F, Q]) and across granular pathological categories.
	\item \textbf{Secondary Classification --- Autonomic Stress Level (Binary):}
	\begin{itemize}
		\item \textit{Class 0 (Resting / Baseline):} Normal parasympathetic dominance.
		\item \textit{Class 1 (Acute Autonomic Stress):} Sympathetic hyperactivation characterized by suppressed RMSSD and elevated LF/HF ratio.
	\end{itemize}
\end{enumerate}

\subsection{In-Silico Dynamical Rhythm Simulation}
To provide predictive insights beyond static scalar measurements, the extracted user PRV features dynamically parameterize a coupled nonlinear Van der Pol oscillator system \cite{vanderpol1928, quiroz2018cardiac}:
\begin{equation}
	\frac{d^2 x}{dt^2} - \mu (1 - x^2)\frac{dx}{dt} + \omega_0^2 x + c_{12}(x - y) = F_{\text{stress}}(t)
\end{equation}
where $x(t)$ represents primary pacemaker dynamics, $\omega_0 = 2\pi(\text{Mean HR}/60)$ represents intrinsic cardiac angular frequency, $\mu$ is the nonlinear damping parameter modulated inversely by vagal RMSSD, $c_{12}$ represents inter-nodal conduction coupling, and $F_{\text{stress}}(t)$ is an autonomic sympathetic driving force parameterized by the measured $\text{LF/HF}$ ratio. Numerical integration via Runge--Kutta RK45 \cite{dormand1980family} simulates limit-cycle phase portraits, visualizing rhythm stability margins and bifurcation breakdown under stress.

\section{Hyperparameter Configurations}

Table~\ref{tab:hyperparameters} summarizes all algorithmic, filter, and windowing hyperparameters established across the system pipeline, detailing their numerical values and engineering rationales.

\begin{table}[!htbp]
	\footnotesize
	\sloppy
	\centering
	\caption{System Hyperparameter Configurations and Rationales}
	\label{tab:hyperparameters}
	\begin{tabular}{|p{2.8cm}|p{1.9cm}|p{2.0cm}|p{5.5cm}|}
		\hline
		\textbf{Parameter} & \textbf{Symbol} & \textbf{Value} & \textbf{Engineering Rationale} \\ \hline
		Accelerometer Sampling Rate & $f_s$ & $100\text{--}200\text{ Hz}$ & Nyquist oversampling for cardiac band ($<40\text{ Hz}$); minimizes IMU battery drain \\ \hline
		Bandpass Lower Cutoff & $f_l$ & $0.5\text{ Hz}$ & Attenuates respiration ($0.1\text{--}0.4\text{ Hz}$) and baseline gravity tilt \\ \hline
		Bandpass Upper Cutoff & $f_h$ & $40.0\text{ Hz}$ & Isolates aortic valve micro-vibrations; cuts high-frequency muscular noise \\ \hline
		Filter Order & $N$ & $4$ & Sharp $24\text{ dB/octave}$ roll-off without numerical instability \\ \hline
		Envelope Window Length & $K$ & $\lfloor 0.2 \cdot f_s \rfloor$ & $200\text{ ms}$ window encompasses the duration of one systolic contraction \\ \hline
		Adaptive Threshold Factor & $\alpha$ & $0.60$ & Balances sensitivity against false positive noise triggers \\ \hline
		Refractory Period & $\delta_{ref}$ & $0.35\text{ s}$ & Prevents double-counting AO peaks; sets upper HR ceiling to $171\text{ bpm}$ \\ \hline
		PPI Outlier Valid Bounds & $[\text{PPI}_{min}, \text{PPI}_{max}]$ & $[333, 1500]\text{ ms}$ & Restricts intervals to physiological cardiac range ($40\text{--}180\text{ bpm}$) \\ \hline
		Tachogram Resampling Rate & $f_{resamp}$ & $4.0\text{ Hz}$ & Standard uniform grid for Welch spectral FFT without aliasing LF/HF bands \\ \hline
		LF Frequency Band & $[\text{LF}_1, \text{LF}_2]$ & $0.04\text{--}0.15\text{ Hz}$ & Baroreceptor reflex and sympathetic Mayer wave power \\ \hline
		HF Frequency Band & $[\text{HF}_1, \text{HF}_2]$ & $0.15\text{--}0.40\text{ Hz}$ & Vagal parasympathetic respiratory sinus arrhythmia power \\ \hline
		Physiological Energy Floor & $\sigma_{min}$ & $0.005\text{ m/s}^2$ & Detects phone left on stationary surface (non-physiological) \\ \hline
		Physiological Energy Ceiling & $\sigma_{max}$ & $0.50\text{ m/s}^2$ & Detects active subject walking/motion (corrupted SCG) \\ \hline
		SNR Acceptance Threshold & $\text{SNR}_{min}$ & $12.0\text{ dB}$ & Minimum quality threshold to accept SCG recording and proceed to classification \\ \hline
	\end{tabular}
\end{table}

\section{Performance Evaluation}

To rigorously assess the performance of the proposed system in alignment with engineering and clinical standards \cite{sokolova2009systematic}, quantitative evaluation protocols are established across all pipeline modules.

\subsection{Evaluation Metrics for Classification (Heart Problem Detection)}
For the supervised machine learning classifiers categorizing cardiac status into \textit{Normal Sinus Rhythm} (Negative) versus \textit{Heart Problem / Arrhythmia} (Positive), performance is quantified using a confusion matrix comprising True Positives ($\text{TP}$), True Negatives ($\text{TN}$), False Positives ($\text{FP}$), and False Negatives ($\text{FN}$):
\begin{enumerate}
	\item \textbf{Accuracy ($\text{Acc}$):} Overall proportion of correctly classified cardiac segments:
	\begin{equation}
		\text{Acc} = \frac{\text{TP} + \text{TN}}{\text{TP} + \text{TN} + \text{FP} + \text{FN}}
	\end{equation}
	\item \textbf{Sensitivity / Recall ($\text{Sens}$):} The ability of the model to correctly identify subjects with actual heart problems or arrhythmias:
	\begin{equation}
		\text{Sens} = \frac{\text{TP}}{\text{TP} + \text{FN}}
	\end{equation}
	\textit{Clinical significance:} High sensitivity is paramount in health screening to minimize False Negatives ($\text{FN}$), ensuring that individuals with arrhythmias are not erroneously classified as healthy. Target: $\text{Sens} \ge 88\%$.
	\item \textbf{Specificity ($\text{Spec}$):} The ability of the model to correctly identify healthy normal sinus rhythms:
	\begin{equation}
		\text{Spec} = \frac{\text{TN}}{\text{TN} + \text{FP}}
	\end{equation}
	\textit{Clinical significance:} High specificity prevents False Positives ($\text{FP}$), avoiding unnecessary user anxiety and clinical overburdening. Target: $\text{Spec} \ge 85\%$.
	\item \textbf{Precision / Positive Predictive Value ($\text{PPV}$):} The proportion of positive arrhythmia predictions that are truly pathological:
	\begin{equation}
		\text{Prec} = \frac{\text{TP}}{\text{TP} + \text{FP}}
	\end{equation}
	\item \textbf{F1-Score ($F_1$):} The harmonic mean of Precision and Sensitivity, providing a balanced evaluation metric under imbalanced class distributions:
	\begin{equation}
		F_1 = 2 \times \frac{\text{Prec} \times \text{Sens}}{\text{Prec} + \text{Sens}}
	\end{equation}
	Target: $F_1 \ge 0.86$.
\end{enumerate}

\subsection{Evaluation Metrics for Numerical Pulse Rate Estimation}
To evaluate the numerical accuracy of systolic peak detection and heart rate extraction against simultaneous reference ECG readings, two standard statistical error metrics are computed across $M$ evaluation windows:
\begin{enumerate}
	\item \textbf{Mean Absolute Error ($\text{MAE}$):}
	\begin{equation}
		\text{MAE} = \frac{1}{M}\sum_{k=1}^M \left| \text{HR}_{\text{est},k} - \text{HR}_{\text{ref},k} \right| \quad (\text{Target: } \text{MAE} < 3.0\text{ bpm})
	\end{equation}
	\item \textbf{Root Mean Squared Error ($\text{RMSE}$):}
	\begin{equation}
		\text{RMSE} = \sqrt{\frac{1}{M}\sum_{k=1}^M (\text{HR}_{\text{est},k} - \text{HR}_{\text{ref},k})^2} \quad (\text{Target: } \text{RMSE} < 4.5\text{ bpm})
	\end{equation}
\end{enumerate}

\subsection{Cross-Validation Strategy}
To prevent data leakage and ensure generalizability across unseen subjects:
\begin{itemize}
	\item Models are evaluated using \textit{Stratified 5-Fold Cross-Validation}, ensuring that proportional distributions of normal and arrhythmic beats are preserved across training and testing partitions.
	\item A patient-stratified split ($80\%$ training, $20\%$ holdout testing) ensures that beats from the same subject do not appear simultaneously in both training and test sets.
	\item Macro-averaged performance metrics are computed across folds alongside confusion matrices to verify consistency across diverse cardiac morphologies.
\end{itemize}
